\exercisesection{Dualizing modules}

\begin{enumerate}
\def\labelenumi{\arabic{enumi})}
\tightlist
\item
  Let A be a local Macaulay ring.
\end{enumerate}

\begin{enumerate}
\def\labelenumi{\alph{enumi})}
\item
  We assume that A admits a dualizing module. \(\Omega\) ; we endow the A-module \(\mathrm{A} \oplus \Omega\) of the A-algebra structure for which \(\Omega\) is an ideal with square zero. Prove that \(\mathrm{A} \oplus \Omega\) is a Gorenstein ring (using a maximal regular sequence, reduce to the case where A is Artinian; observe that the socle of \(\mathrm{A} \oplus \Omega\) is identified with that of the A-module \(\Omega\) ).
\item
  Deduce that for A to admit a dualizing module, it is necessary and sufficient that there exist a Gorenstein ring B and a surjective homomorphism from B to A.
\end{enumerate}

\begin{enumerate}
\def\labelenumi{\arabic{enumi})}
\setcounter{enumi}{1}
\tightlist
\item
  Let A be a regular local ring, I an ideal of A, B the ring A/I; suppose that B is a Macaulay ring. Let (L,p) be a minimal projective resolution of the A-module B; it has length \(c = \text{ht}(I)\) .
\end{enumerate}

\begin{enumerate}
\def\labelenumi{\alph{enumi})}
\item
  Let \(\mathrm{L}^*\) the complex \(\operatorname{Homgr}_{\mathrm{A}}(\mathrm{L},\mathrm{A})\) ; prove that we have \(\mathrm{H}^{i}(\mathrm{L}^{*}) = 0\) for \(i\neq c\) and that the B-module \(\mathrm{H}^c (\mathrm{L}^*)\) is dualizing.
\item
  Prove that the following conditions are equivalent:
\end{enumerate}

\begin{enumerate}
\def\labelenumi{(\roman{enumi})}
\item
  A is a Gorenstein ring;
\item
  the complexes L and \(L^{*}(-d)\) are isomorphic;
\item
  \(L_{c}\) is free of rank one.
\end{enumerate}

\begin{enumerate}
\def\labelenumi{\arabic{enumi})}
\setcounter{enumi}{2}
\tightlist
\item
  Let A be a ring admitting a dualizing module \(\Omega\) .
\end{enumerate}

\begin{enumerate}
\def\labelenumi{\alph{enumi})}
\item
  For \(\Omega\) is isomorphic to an ideal I of A, it is necessary and sufficient that \(A_{p}\) let R be a Gorenstein ring for every minimal prime ideal \(p\) of A (observe that \(\Omega\) is isomorphic to a submodule of \(\bigoplus \Omega_{p}\) ).
\item
  When these conditions are satisfied, prove that I is equal to A or has height 1, and that A/I is a Gorenstein ring (compute the depth of A/I, then the A/I-module \(\mathrm{Ext}_{\mathrm{A}}^{1}(\mathrm{A / I,I})\) ).
\item
  Prove that a factorial ring admitting a dualizing module is a Gorenstein ring.
\end{enumerate}

\begin{enumerate}
\def\labelenumi{\arabic{enumi})}
\setcounter{enumi}{3}
\item
  Let A be a local Macaulay ring of dimension \(d\) ; let \(\Omega\) a finitely generated A-module, of finite injective dimension, of depth \(d\) , such that the canonical homomorphism \(\mathrm{A} \to \mathrm{End}_{\mathrm{A}}(\Omega)\) is bijective. Prove that the \(\Lambda\) -module \(\Omega\) is dualizing (argue by induction on \(d\) , deducing from exerc. 7 of § 3 that for every element \(\Omega\) -regular \(x\) of \(\mathfrak{m}_{\mathrm{A}}\) , the canonical homomorphism \(\mathrm{A}/x\mathrm{A} \to \mathrm{End}_{\mathrm{A}/x\mathrm{A}}(\Omega/x\Omega)\) is bijective).
\item
  Let A be a Noetherian local ring admitting a dualizing module \(\Omega\) . Let T be a finitely generated A-module of finite injective dimension.
\end{enumerate}

\begin{enumerate}
\def\labelenumi{\alph{enumi})}
\item
  Prove that \(\mathrm{Ext}_{\mathrm{A}}^{i}(\Omega, \mathrm{T})\) is zero for \(i > 0\) (use exerc. 7 of § 3).
\item
  Prove that the canonical homomorphism \(\Omega\otimes_{\Lambda}\mathrm{Hom}_{\mathbb{A}}(\Omega,\mathrm{T})\to\mathrm{T}\) is bijective (apply a) and exerc. 10 of § 8).
\item
  Prove using exercise 8, b) of § 8 the equalities dp \(_{A}\) Hom \(_{A}\) (Ω,T)) = dim(A) - prof(T) and prof(Hom \(_{A}\) (Ω,T)) = prof(T). Deduce that if prof(T) = dim(A), then the A-module T is a direct sum of modules isomorphic to Ω.
\end{enumerate}

\(d)\) We set \(c = \dim(A) - \operatorname{prof}(T)\) Prove that there exist integers \(n_0, \ldots, n_c\) and an exact sequence

\[
0 \rightarrow \Omega^ {n _ {c}} \rightarrow \dots \rightarrow \Omega^ {n _ {0}} \rightarrow \mathrm{T} \rightarrow 0
\]

(argue by induction on c, constructing with the aid of b) a surjective homomorphism \(\Omega^{n_{0}} \to T\) ).

¶ 6) Let A be a local noetherian ring admitting a dualizing module Ω; let M be a finitely generated A-module of finite projective dimension.

\begin{enumerate}
\def\labelenumi{\alph{enumi})}
\item
  Prove that one has \(\operatorname{Tor}_{i}^{\mathrm{A}}(\Omega, \mathrm{M}) = 0\) for \(i > 0\) (Argue by induction on \(\dim(\mathrm{A})\) , considering an exact sequence \(0 \to \mathrm{N} \to \mathrm{L} \to \mathrm{M} \to 0\) where \(\mathrm{L}\) is free of finite type. If \(x\) is a non-zero-divisor element of \(\mathfrak{m}_{\mathrm{A}}\) , deduce from the induction hypothesis that \(\operatorname{Tor}_{i}^{\mathrm{A}}(\Omega, \mathrm{N})\) is zero for \(i > 0\) , which implies \(\operatorname{Tor}_{i}^{\mathrm{A}}(\Omega, \mathrm{M}) = 0\) for \(i > 1\) , and that the homothety of ratio \(x\) is injective in \(\Omega \otimes_{\mathrm{A}} \mathrm{N}\) , which implies that every prime ideal \(\mathfrak{p}\) associated with \(\operatorname{Tor}_{1}^{\mathrm{A}}(\Omega, \mathrm{M})\) is a minimal ideal of \(\operatorname{Spec}(\mathrm{A})\) ; observe that at such a \(\mathfrak{p}\) the \(\mathrm{A}_{\mathfrak{p}}\) -module \(\mathrm{M}_{\mathfrak{p}}\) is free).
\item
  Prove that the A-module \(\Omega\otimes_{\mathrm{A}}\mathrm{M}\) is of finite injective dimension.
\item
  Prove that the map \(\theta: \mathrm{M} \to \mathrm{Hom}_{\Lambda}(\Omega, \Omega \otimes_{\Lambda} \mathrm{M})\) defined by \(\theta(m)(w) = w \otimes m\) for \(m \in \mathrm{M}\) , \(w \in \Omega\) is an isomorphism (reduce with the aid of a projective resolution to the case where M is free).
\end{enumerate}

\(d\) ) The maps \([\mathrm{T}] \mapsto [\mathrm{Hom}_{\mathrm{A}}(\Omega, \mathrm{T})]\) and \([\mathrm{M}] \mapsto [\Omega \otimes_{\mathrm{A}} \mathrm{M}]\) define mutually inverse bijections between the set of isomorphism classes of finitely generated A-modules of finite injective dimension and the set of isomorphism classes of finitely generated A-modules of finite projective dimension, * and quasi-inverse equivalences of categories between the corresponding categories*.

¶ 7) Let A be a macaulay local ring, and M a finitely generated macaulay A-module. We denote by \(r(M)\) (resp. \(t(M)\) ) the dimension of the \(\kappa_{A}\) -vector space \(\kappa_{A} \otimes_{A} M\) (resp. \(\operatorname{Ext}_{A}^{p}(\kappa_{A}, M)\) , with \(p = \operatorname{prof}(M)\) ), cf. exercise 16 of § 1.

\begin{enumerate}
\def\labelenumi{\alph{enumi})}
\item
  Suppose that A possesses a dualizing module \(\Omega\) ; set \(M' = \operatorname{Ext}_A^c(M, \Omega)\) , with \(c = \dim(A) - \dim_A(M)\) . Prove the equalities \(r(M') = t(M)\) and \(t(M') = r(M)\) (reduce with the aid of prop. 7 of § 3 to the case \(c = 0\) , then by passing to the quotient by a maximal secant sequence to the case where A is artinian, and use prop. 6 of § 8).
\item
  Hence we have \(t(\mathrm{M}_{\mathfrak{p}}) \leqslant t(\mathrm{M})\) for every \(\mathfrak{p} \in \operatorname{Spec}(\mathrm{A})\) (reduce with the aid of exercises 2 of § 2 and 16 of § 1 to the case where A is complete, and apply a)).
\item
  For the module M to be dualizing, it is necessary and sufficient that it be faithful and that one have \(t(\mathrm{M}) = 1\) (reduce to the case where A is complete; apply a) and the corollary of th. 1).
\end{enumerate}

\begin{enumerate}
\def\labelenumi{\arabic{enumi})}
\setcounter{enumi}{7}
\tightlist
\item
  Let A be a local Gorenstein ring, a a non-zero ideal of height zero, and b the annihilator of a.
\end{enumerate}

\begin{enumerate}
\def\labelenumi{\alph{enumi})}
\item
  Prove that the ideal b is of height zero, and that the A-module A/b has no embedded associated prime ideals (observe that every prime ideal associated to A/b is associated to A).
\item
  Suppose that the A-module A/α has no embedded associated prime ideals; prove that α is the annihilator of β (reduce to the case where dim(A) = 0, and use Matlis duality).
\item
  For A/a to be a Macaulay ring, it is necessary and sufficient that the same be true of A/b (use the corollary of th. 1 and that of prop. 3). If these conditions are satisfied, the A/a-module b is dualizing.
\end{enumerate}

\begin{enumerate}
\def\labelenumi{\arabic{enumi})}
\setcounter{enumi}{8}
\tightlist
\item
  Let \(\Gamma\) a submonoid of the additive monoid \(\mathbf{N}\) such that \(\mathbf{N} - \Gamma\) be finite; we denote by \(c\) the smallest integer such that \(\Gamma\) contains all the integers \(\geqslant c\) . Let \(k\) a field, B the ring \(k[[T]]\) , \(\mathrm{K} = k((\mathrm{T}))\) its field of fractions, and A the sub- \(k\) -algebra of B generated by the elements \(\mathrm{T}^{\gamma}\) for \(\gamma \in \Gamma\) .
\end{enumerate}

\begin{enumerate}
\def\labelenumi{\alph{enumi})}
\item
  Prove that A is a local integral noetherian ring of dimension 1, whose field of fractions is K and whose integral closure is B; the ideal c = A : B is equal to BT \(^{c}\) .
\item
  For A to be a Gorenstein ring, it is necessary and sufficient that one have \(c = 2\operatorname{Card}(\mathbf{N} - \Gamma)\) .
\item
  Let \(\Omega\) the sub-A-module of K generated by the elements \(T^{-\alpha}\) for \(\alpha \notin \Gamma\) . Prove that \(K/\Omega\) is a Matlis A-module (use prop. 2 of § 8), and that \(\Omega\) is a dualizing A-module.
\end{enumerate}

\begin{enumerate}
\def\labelenumi{\arabic{enumi})}
\setcounter{enumi}{9}
\tightlist
\item
  Let A be a noetherian ring. Let I be a bounded injective A-complex, whose homology is of finite type; for every A-complex C, we denote by D(C) the complex Homgr \(_{A}\) (C,I).
\end{enumerate}

\begin{enumerate}
\def\labelenumi{\alph{enumi})}
\tightlist
\item
  Prove that the following conditions are equivalent:
\end{enumerate}

\begin{enumerate}
\def\labelenumi{(\roman{enumi})}
\item
  for every finitely generated A-module M, the canonical morphism of complexes \(\boldsymbol{\alpha}_{\mathrm{M}}:\mathrm{M}\to\mathbf{D}(\mathbf{D}(\mathrm{M}))\) (n° 5) is a homomorphism;
\item
  for every bounded A-complex C whose homology is of finite type, the morphism \(\alpha_{\mathrm{C}}: \mathrm{C} \to \mathbf{D}(\mathbf{D}(\mathrm{C}))\) (loc. cit.) is a homologism;
\item
  the canonical morphism \(\alpha_{\mathrm{A}}:\mathrm{A}\to \mathrm{Homgr}_{\mathrm{A}}(\mathrm{I},\mathrm{I})\) is a homological isomorphism.
\end{enumerate}

(Adapt the proof of Thm. 1.)

We say that the complex I is dualizing if it is bounded, injective, H(I) is finitely generated, and it satisfies the equivalent conditions above.

\begin{enumerate}
\def\labelenumi{\alph{enumi})}
\setcounter{enumi}{1}
\item
  If A admits a dualizing module \(\Omega\) any finite-length injective resolution (I, δ) of \(\Omega\) defines a dualizing complex (loc. cit.).
\item
  Let B be an A-algebra which is a finitely generated A-module, and let I be a dualizing complex of A-modules; prove that the complex of B-modules \(\mathrm{Homgr}_{\mathrm{A}}(\mathrm{B},\mathrm{I})\) is dualizing. Thus every quotient ring of a Gorenstein ring of finite dimension admits a dualizing complex.
\item
  Let I be a bounded injective A-complex. For I to be dualizing, it is necessary and sufficient that the complex of \(A_{m}\) -modules \(I_{m}\) so that it is dualizing for every maximal ideal m of A.
\item
  Let I be a dualizing A-complex, P a projective A-module of rank 1, n an integer. The complex \(1 \otimes_{\mathrm{A}} \mathrm{P}(n)\) is dualizing.
\item
  Let I, J be bounded injective A-complexes whose homology is of finite type, and \(u: I \rightarrow J\) a homomorphism; for J to be dualizing, it is necessary and sufficient that I be so.
\end{enumerate}

¶ 11) Let A be a Noetherian ring.

\begin{enumerate}
\def\labelenumi{\alph{enumi})}
\tightlist
\item
  Assume \(\Lambda\) local. Let P, Q be bounded flat A-complexes. Suppose that \(\mathrm{H}(\mathrm{P}\otimes_{\mathrm{A}}\mathrm{Q})\) is zero in degree \(\neq 0\) and free of rank 1 in degree 0. Prove that there exists an integer \(p\in \mathbf{Z}\) and homologisms \(\mathrm{A}\to \mathrm{P}(p)\) and \(\mathrm{A}\to \mathrm{Q}(-p)\) .
\end{enumerate}

(Using lemma 1 and prop. 1 of A, X, pp.~66 and 62, reduce to the case where P and Q are zero on the right; using prop. 1 of loc. cit. and th. 3 of II, § 5, no. 4, then construct complexes P' and Q' such that P = A ⊕ P' and Q = A ⊕ Q', and prove that one has H(P') = H(Q') = 0.)

\begin{enumerate}
\def\labelenumi{\alph{enumi})}
\setcounter{enumi}{1}
\item
  Let I, J be dualizing A-complexes (exerc. 10). If A is local, prove that there exists an integer n and a homologism of J onto I(n) (let P = Homgr \(_A\) (I, J), Q = Homgr \(_A\) (J, I); using the morphism \(\nu\) of no. 7 and the homologism \(\alpha_{J}\) : J → Homgr \(_A\) (Q, I) relative to the dualizing complex 1, construct a homomorphism P ⊗ \(_A\) Q → A, and apply a)).
\item
  In the general case, prove that there exists an integer n, a projective A-module L of rank 1, and a homomorphism from J onto \(I \otimes_{A} L(n)\) (set \(L = H(\text{Homgr}_{A}(I, J))\) , and apply b)).
\end{enumerate}

\begin{enumerate}
\def\labelenumi{\arabic{enumi})}
\setcounter{enumi}{11}
\tightlist
\item
  Let \(k\) be a field, and R a \(k\) -graded algebra of type \(\mathbf{N}\) such that \(\mathrm{R}_0 = k\) and that R be a Macaulay ring, of dimension \(d\) For every finitely generated graded R-module M, we denote by \(\mathrm{P}_{\mathrm{M}}(\mathrm{T})\) the Poincaré series \(\sum_{i} (\dim_k \mathrm{M}_i) \mathrm{T}^i\) and one sets \(\mathrm{Q}_{\mathrm{M}}(\mathrm{T}) = (1 - \mathrm{T})^{\dim(M)} \mathrm{P}_{\mathrm{M}}(\mathrm{T})\) ; it is an element of \(\mathbf{Z}[\mathrm{T}, \mathrm{T}^{-1}]\) (VIII, § 6, no. 3, prop. 5).
\end{enumerate}

\begin{enumerate}
\def\labelenumi{\alph{enumi})}
\item
  Let \(\Omega\) a dualizing R-module; prove that \(\Omega\) admits a graded R-module structure for which \(Q_{\Omega}(T) = (-1)^{d} Q_{\mathbb{R}}(T^{-1})\) (write R as a quotient of a polynomial algebra \(A = k[X_1, \ldots, X_n]\) , with \(\deg(X_i) = d_i\) ; if L is a free graded resolution of the A-module R, observe that the complex \(\mathrm{Homgr}_{\mathrm{A}}(\mathrm{L}, \mathrm{A})(-d)\) defines one of \(\Omega\) , and use Example 3 of VIII, § 4, no. 2).
\item
  Assume that R is a Gorenstein ring; prove that there exists an integer \(a \in \mathbf{Z}\) satisfying \(\mathrm{Q}_{\mathbb{R}}(\mathrm{T}^{-1}) = (-1)^{d}\mathrm{T}^{a}\mathrm{Q}_{\mathbb{R}}(\mathrm{T})\) .
\item
  Assume that R is an integral domain and that there exists an integer \(a \in \mathbf{Z}\) such that one has \(\mathrm{Q}_{\mathrm{R}}(\mathrm{T}^{-1}) = (-1)^{d}\mathrm{T}^{a}\mathrm{Q}_{\mathrm{R}}(\mathrm{T})\) Prove that R is a Gorenstein ring (endow \(\Omega\) of a graduation for which \(\mathrm{Q}_{\Omega} = \mathrm{Q}_{\mathrm{R}}\) , and prove that a nonzero element of \(\Omega_{0}\) forms a basis of \(\Omega\) ).
\item
  Let R be the graded k-algebra \(k[X,Y]/(X^{3},XY,Y^{2})\) ; prove that we have \(Q_{R}(T^{-1}) = T^{-2}Q_{R}(T')\) , but that R is not a Gorenstein ring.
\end{enumerate}

